\section{Verification, reproducibility, and limits}
\label{sec:validation}

The infinite mathematical assertions are established by the proofs in
this report. The accompanying programs check the finite constructions,
exact constants, and arithmetic certificates used by those proofs.
Their successful output is not a substitute for a proof of an
asymptotic density statement, a Galois-cohomology comparison, or a new
zero-location theorem.

\subsection{Frozen sources and the earlier-report corpus}

The source inspection uses \texttt{lsernau/grh-lean} at commit
\begin{center}\small
\texttt{72fead102f663d5c7a407964a13dc0cf9d977207}.
\end{center}
The archive contains 31 selected upstream source and documentation
files, including 25 Lean files. Their bytes were checked against the
Git blob SHA-1 identifiers returned for this commit, and SHA-256
digests and byte counts are recorded in
\path{sources/upstream_manifest.json}. The original license, notice,
and attribution files are retained. These checks establish the
identity of the inspected text.

The ProveIt corpus is frozen at commit
\texttt{5cc267af75a635199ac5d93a7d49cbce521daefc}.
All 46 relevant incoming ZIP archives, totalling 43,822,390 bytes,
were downloaded and checked against their Git blob identifiers;
their SHA-256 digests and sizes are recorded in
\path{data/archive_verification.json}. The inventory is in
\path{data/archive_inventory.csv}. Distinct filenames, including
numbered variants, were kept distinct. The review used the extracted
reports to identify previous results, correct the Selmer-precision
claim, and avoid presenting inherited finite computations as new.
Hash verification does not validate the mathematics inside an
earlier archive, and this review is not an independent reproof of
every claim in those reports.

The upstream audit is static. No complete build of either Lean
edition and no independent kernel replay were run in this
environment; Lean and Lake were not installed. The upstream
README reports its own build, axiom, and independent-checker results.
Those reports are distinguished from checks performed here.
The new combined statements in this report have human-readable
proofs; they are not described as newly Lean-checked.
The exact source anchors and audit limitations are retained in
\path{sources/audit_findings.json} and
\path{sources/upstream_audit.md}.

\subsection{Finite certificates and their scope}

\paragraph{Ideal-density constructions.}
The program \path{code/check_ideal_density.py} checks the
five-point exponent exchange, 10,000 ternary saturation instances,
five finite smooth-monoid examples, and the cube exchange in
dimensions \(2,3,4\). It also verifies explicit signed-step
counterexamples showing why the cube-exchange construction is not
valid for arbitrary fractional-ideal ratios.

For \(K=\mathbb Q(\sqrt{-23})\), the density-gap certificate uses
\(Y=10^7\). The prime-ideal product has 664,449 factors: two factors
for each of 332,112 split rational primes, 224 inert-prime-ideal
factors of norm \(p^2\), and one ramified factor. The product is
rounded downward after every multiplication using an integer
denominator \(10^{40}\), so the reported value is a rigorous lower
bound rather than a floating-point estimate. With the proved tail
bound \(2/Y\), the recorded rational gap exceeds \(1/20{,}000{,}000\).
The full output is \path{data/ideal_density_certificate.json}.

The separate program \path{code/verify_ideal_natural_gain.py}
independently verifies the rational arithmetic for the four selected
prime ideals of norms \(2,2,3,3\), the simple splitting of \(2\) and
\(3\) in the defining quadratic order, and the elementary zeta-value
upper bounds used in the proof. It certifies the strict natural-density
gain
\[
 \frac{16658080679}{13931406950400}>0.0011957.
\]
Its output is \path{data/ideal_natural_gain_certificate.json}.
This particular gain concerns proper integral-ideal ratios.
The all-number-fields extension and the natural-versus-upper-density
separation follow from the general arguments in the manuscript,
not from checking finitely many fields.

\paragraph{Auxiliary-character conductor minima.}
The program \path{code/check_source_twists.py} checks five minima
under the source-compatible precursor condition and six minima
under the weaker condition on the final anticyclotomic character.
It uses exact modular arithmetic and arithmetic-progression
Eratosthenes sieves. The progression sieve is cross-checked against
an ordinary sieve through \(10^4\); modular square roots are
cross-checked by exhaustive search at odd primes through \(500\).

The certificate includes the exclusion of every smaller permitted
multipair support using the reductions proved in
Section~\ref{sec:source_twists}. In the exceptional two-pair search,
five support pairs are tested with both orientations, and the
corresponding relation determinants are nonzero modulo \(3\).
The outputs are
\path{data/source_twist_certificate.json} and
\path{data/source_twist_minima.csv}.
The class-field interpretation, the scope of the admissible
characters, and the reduction to these finite searches are proved
in the text. The calculation does not establish the required
automorphic determinant estimates or a nonabelian RH.

\paragraph{Module-theoretic interfaces.}
The program \path{code/check_module_bridge.py} performs 46,332
one-level saturation checks and 50,193 finite-certificate equivalence
checks on explicitly specified torsion-module types. It also
enumerates 30 finite cyclic-cocycle scalar-extension models,
including oddness in residue characteristic \(2\).
The results are in \path{data/module_bridge_checks.json}.
These are exact finite sanity checks for
Section~\ref{sec:sha-odd-bridge}; they do not compute the arithmetic
groups \(\mathcal A_{\ell,j,\nu}(E;\mathbb Z_\ell)\).
The proof of scalar extension and the compactness argument are
supplied separately in that section.

\paragraph{Inherited \(K\)-theory preparation.}
The prior-report verifier \path{code/verify_ktheory_prefilter.py}
was replayed with exact rational arithmetic. For the 14 even
parameters \(54\leq e\leq80\), it compares an exponential-series
recurrence with the Bernoulli-polynomial formula in 84 evaluations,
checks the corresponding greatest common divisors and explicit
B\'ezout certificates, and computes eight compact invariant-space
dimensions. The record is
\path{data/ktheory_prefilter_rerun.json}.
In particular, the two displayed gcds for degrees \(124\) and
\(128\) are \(4\). These computations are inherited preparatory
certificates. They neither compute the missing Hecke determinants
nor certify a new algebraic \(K\)-group vanishing.

\subsection{Replaying the supplied computations}

The five programs use Python's standard library only. Run the
following commands, with assertions enabled, from the directory
containing \path{report.tex}; the supplied \path{data/} directory
receives the records. Python 3.10 or later is sufficient.

\begin{minipage}{\textwidth}
\begin{verbatim}
python3 code/check_ideal_density.py --cutoff 10000000 \
  --output data/ideal_density_certificate.json

python3 code/verify_ideal_natural_gain.py \
  --output data/ideal_natural_gain_certificate.json

python3 code/check_source_twists.py

python3 code/check_module_bridge.py \
  --output data/module_bridge_checks.json

python3 code/verify_ktheory_prefilter.py --ranks \
  --output data/ktheory_prefilter_rerun.json
\end{verbatim}
\end{minipage}

The character verifier writes both its JSON certificate and its CSV
table automatically. Verification logs are included under
\path{verification/}. Replaying these commands reproduces the
finite certificates; rebuilding the upstream Lean development
requires its separately pinned Lean and Mathlib environment.

The resulting evidence supports the distinctions used throughout
this report: proved general statements, exact finite certificates,
source theorem extractions, and conditional arithmetic transfers.
No finite prime search is used as evidence that all zeros lie on a
line, and no small normalized length is identified with vanishing
without the hypotheses that justify that inference.

